\section{曲率区分凸、严格凸与强凸}
\label{sec:09-Convex-Optimization/03-Convex-Functions/03}

一个线性模型，两个特征几乎共线（比如``月消费额''和``季度消费额除以三''）。换个求解器重跑一遍，第一次得到的系数是 \((3.1,-2.0)\)，第二次是 \((0.4,0.7)\)，训练损失都是 \(0.2170\)，小数点后四位一模一样。业务方拿着第一组系数问：为什么月消费越高，风险评分越低？这个问题没法回答，因为第二组系数会给出相反的故事，而两组在数据上同样好。

加一项 \(\lambda\norm{x}_2^2\)，\(\lambda=10^{-3}\)，重跑。两次结果一致到小数点后六位。

这里没有出现过任何非凸的东西。前后两个目标都是凸的，问题只在于第一个目标的谷底是平的：沿着 \(x_1+x_2\) 不变的那条直线走，损失完全不动，最优解是一整条线段而不是一个点。凸性从来没有承诺过唯一性，它只承诺弦在图像上方。

平坦还有第二个症状。\(f(x)=x^2\) 上跑梯度下降几步到位，\(f(x)=x^4\) 上磨磨蹭蹭，越靠近原点越慢。两个都是凸的，\(x^4\) 甚至连谷底都不平——它的最低点唯一。差别在于底部附近的曲率：\(x^2\) 的二阶导恒为 \(2\)，\(x^4\) 的二阶导 \(12x^2\) 在原点处消失。

\begin{center}
\begin{tikzpicture}[x=0.62cm,y=1.25cm,font=\small,>=stealth]
  \draw[->,black!55] (-2.35,0) -- (2.45,0);
  \draw[blue!65!black,very thick] plot[domain=-2.05:-1,samples=30] (\x,{(\x+1)*(\x+1)+0.35});
  \draw[blue!65!black,very thick] plot[domain=1:2.05,samples=30] (\x,{(\x-1)*(\x-1)+0.35});
  \draw[orange!85!black,line width=2.2pt] (-1,0.35) -- (1,0.35);
  \node[font=\footnotesize,align=center,black!75] at (0,-0.45) {最优解是一整段};
  \node[font=\footnotesize,blue!65!black] at (0,2.05) {凸};
  \begin{scope}[shift={(5.2,0)}]
    \draw[->,black!55] (-2.35,0) -- (2.45,0);
    \draw[gray!75,dashed,thick] plot[domain=-2.05:2.05,samples=50] (\x,{0.35+0.15*\x*\x});
    \draw[blue!65!black,very thick] plot[domain=-2.05:2.05,samples=70] (\x,{0.35+0.08*\x*\x*\x*\x});
    \fill[blue!65!black] (0,0.35) circle (2.2pt);
    \node[font=\footnotesize,align=center,black!75] at (0,-0.45) {最低点唯一，底部趋平};
    \node[font=\footnotesize,blue!65!black] at (0,2.05) {严格凸};
  \end{scope}
  \begin{scope}[shift={(10.4,0)}]
    \draw[->,black!55] (-2.35,0) -- (2.45,0);
    \draw[gray!75,dashed,thick] plot[domain=-2.05:2.05,samples=50] (\x,{0.35+0.15*\x*\x});
    \draw[blue!65!black,very thick] plot[domain=-2.05:2.05,samples=60] (\x,{0.35+0.32*\x*\x});
    \fill[blue!65!black] (0,0.35) circle (2.2pt);
    \node[font=\footnotesize,align=center,black!75] at (0,-0.45) {抛物线始终托在下方};
    \node[font=\footnotesize,blue!65!black] at (0,2.05) {强凸};
  \end{scope}
\end{tikzpicture}

\emph{三幅图里的虚线抛物线是同一条：中间那幅它戳穿了曲线，右边那幅它稳稳托在下面，差别就是强凸与否。}
\end{center}

三个假设一级比一级强。凸保证不会更差，严格凸再买到唯一性，强凸买到的是速率。买到什么、付出什么，下面逐级对账。

\subsection{二阶判据和一阶下界是同一件事的两个面}

先把曲率和上一篇的切平面下界接上。设 \(f\) 在开凸域上二次可微，则

\[
f\ \text{凸}
\iff
\nabla^2f(x)\succeq0\ \text{对所有 }x.
\]

半正定的含义是：对任意方向 \(v\)，\(v^{\trans}\nabla^2f(x)v\ge0\)。把 \(f\) 限制在过 \(x\)、方向为 \(v\) 的直线上，令 \(\phi(t)=f(x+tv)\)，这个二次型正是 \(\phi''(0)\)。所以高维凸性可以按方向逐条读：不管朝哪走，脚下的曲面都不向下弯。

从曲率到下界只需要一次带积分余项的 Taylor 展开。取 \(d=y-x\)，

\[
f(y)=f(x)+\nabla f(x)^{\trans}d+\int_0^1(1-t)\,d^{\trans}\nabla^2f(x+td)\,d\;\mathrm{d}t.
\]

Hessian 处处半正定时被积函数非负（\(1-t\ge0\)），积分项非负，直接得到 \(f(y)\ge f(x)+\nabla f(x)^{\trans}(y-x)\)。反向由 \(\phi''(0)\ge0\) 逐点读回来。

条件对一遍。这条链子要求二次可微，比一阶下界多了一层假设，因此它是判据而不是定义——\(\abs{x}\) 是凸的，但 Hessian 在原点不存在。积分路径 \(x+td\) 必须整条落在定义域内，这由定义域是凸集保证。还要注意积分沿的是连接两点的线段，所以``只在某一个点验过 Hessian 半正定''推不出任何全局结论，二阶条件必须在整个凸域上逐点成立。

\subsection{三个定义，三级递增的要求}

普通凸性允许平坦。仿射函数既凸又凹，弦不等式处处取等号，梯度下降原地不动也合法，因为处处都是最优。开头那个共线特征的例子就落在这一级：Hessian 是 \(X^{\trans}X\)，共线让它有一个近乎零的特征值，对应方向上目标几乎不变。

严格凸掐掉取等的情形：对 \(x\ne y\) 和 \(0<\theta<1\)，

\[
f(\theta x+(1-\theta)y)<\theta f(x)+(1-\theta)f(y).
\]

推论是唯一性。若 \(x^\star\ne y^\star\) 都是最优点，取它们的中点，严格不等式给出比最优值更低的目标值，矛盾。注意这句话的前提是最优解存在——\(e^x\) 严格凸，但在 \(\R\) 上取不到最小值。很多收敛证明把唯一性当作隐含前提，出问题时往往就出在这里。

强凸再加一个定量的托底：存在 \(m>0\)，使对所有 \(x,y\)，

\[
f(y)\ge f(x)+\nabla f(x)^{\trans}(y-x)+\frac m2\norm{y-x}_2^2.
\]

图像不止在切平面上方，还高出一个曲率为 \(m\) 的二次碗。二次可微时它等价于 \(\nabla^2f(x)\succeq mI\)：每个方向、每个位置的曲率都至少是 \(m\)。图里右边那幅的虚线抛物线就是这个碗，它在整条曲线下方。

一个等价说法有时更好用：\(f\) 是 \(m\)-强凸，当且仅当 \(f(x)-\frac m2\norm{x}_2^2\) 仍然是凸函数。这个形式不需要可微性，因而适用于非光滑的目标。

\subsection{\texorpdfstring{\(x^4\) 把严格与强的差别摆在明处}{x^4 把严格与强的差别摆在明处}}

\(f(x)=x^4\) 在实轴上严格凸，最低点唯一。但 \(f''(x)=12x^2\) 在原点为零，因此不管把 \(m\) 取得多小，总能找到足够靠近原点的 \(x\) 使 \(f''(x)<m\)。全局统一的曲率下界不存在，它严格凸而不强凸。图里中间那幅画的就是这件事：虚线抛物线在原点附近戳到了曲线上方，说明它不是一个合法的下界。

\(f(x)=x^2\) 的二阶导恒为 \(2\)，是 \(2\)-强凸。两个函数的最优解相同、唯一性相同，差别只在曲率的全局下界，而这个下界决定了算法要花多少步。

由此得到一个用得非常频繁的配方：给任意凸函数加上 \(\frac\lambda2\norm{x}_2^2\)，\(\lambda>0\)，和函数至少是 \(\lambda\)-强凸。用上面那个等价说法一句话就能看出来：和减去 \(\frac\lambda2\norm{x}_2^2\) 剩下原函数，还是凸的。开头那个岭回归的修复就是这么起作用的——它把 \(X^{\trans}X\) 的谱整体抬高 \(\lambda\)，那个近乎零的特征值变成 \(\lambda\)，平坦的谷底被顶成一个有唯一底点的碗。系数从此稳定，代价是引入了偏差，这笔账在\hyperref[sec:09-Convex-Optimization/07-Applications/01]{``光滑经验风险与 \(L_2\) 正则化''}里算。

\subsection{强凸把目标值的差翻译成距离}

在最优点 \(x^\star\) 处 \(\nabla f(x^\star)=0\)，代进强凸定义（取 \(x=x^\star\)、\(y=x\)）得

\[
f(x)-f(x^\star)\ge\frac m2\norm{x-x^\star}_2^2.
\]

这一步把``目标值接近最优''换算成了``参数接近最优参数''。它值钱的地方在于：算法通常只能观测目标值和梯度，而人关心的往往是参数本身。普通凸函数上这个换算不成立——谷底是平的时候，可以在谷底两端来回晃，目标值分毫不差，参数却相距很远，开头那两组系数正是如此。

反方向的假设是曲率上界。梯度 \(L\)-Lipschitz 在二次可微时等价于 \(\nabla^2f(x)\preceq LI\)。\(m\) 卡住最小曲率，\(L\) 卡住最大曲率，比值

\[
\kappa=\frac Lm
\]

是这个优化问题的条件数。几何上想象一个又长又窄的椭圆碗：短轴方向曲率大、梯度陡，长轴方向曲率小、梯度缓。固定步长的梯度下降必须迁就短轴（否则发散），于是在长轴上爬得很慢，误差按 \((1-1/\kappa)^k\) 的速度线性下降。\(\kappa\) 直接出现在步数里，缩放、预条件和 Newton 法都是在设法把椭圆碗压圆，见\hyperref[sec:09-Convex-Optimization/06-Optimization-Algorithms/01]{``梯度下降真正难在步长''}与\hyperref[sec:09-Convex-Optimization/06-Optimization-Algorithms/02]{``Newton 法用曲率校正方向''}。

没有强凸就没有这个线性速率。只有凸加 \(L\)-光滑时，梯度下降的目标值误差按 \(O(1/k)\) 走，慢一个量级；只有凸加次梯度界时更慢，\(O(1/\sqrt k)\)。三级假设换来的三档速率，是选择正则化强度时实际要权衡的东西：\(\lambda\) 调大，\(\kappa\) 变小、优化更快，同时偏差变大、拟合更差。

\subsection{Hessian 正定是充分条件，不是必要条件}

几个容易踩的坑值得单独点一下。

Hessian 处处正定推出严格凸，反过来不成立：\(x^4\) 在原点的二阶导是零，它照样严格凸。``处处正定''是严格凸的充分条件。

强凸不需要可微。一个强凸的光滑函数加上任意凸的非光滑函数，和仍然强凸，因为定义只要求那个二次碗托在下面。弹性网 \(\lambda_1\norm{x}_1+\frac{\lambda_2}2\norm{x}_2^2\) 加在任何凸损失上都得到强凸目标，尽管它在坐标轴上处处不可微。Hessian 是方便的入口，不是唯一入口。

强凸也不等于条件数好。\(m=10^{-6}\) 时函数确实强凸，所有定理都适用，但 \(\kappa\) 可能是 \(10^{9}\)，速率常数大到没有实际意义。定性的``满足假设''和定量的``跑得快''是两件事，报告收敛保证时要把 \(m\) 和 \(L\) 的实际数值一并说出来。

判断顺序可以固定下来：先确认凸（验弦不等式或查 Hessian 半正定），再看能否排除平坦方向以确认严格凸，最后找有没有全局统一的 \(m>0\) 以确认强凸。每一级要求更高，也换回更强的保证。

麻烦在于，真实目标函数很少长得像 \(x^2\) 或 \(x^4\)。一个带 softmax、\(\log\)、\(\max\) 和若干矩阵乘法的损失，写下来占半页纸，几万维参数的 Hessian 既算不动也存不下。下一篇换一条路：不去验证结果，去检查它是怎么被造出来的。
